% Standalone problem document, generated from the adjacent README.md.
% Compile with XeLaTeX. Mathematical content is maintained in Markdown.
\documentclass[11pt,a4paper]{article}
\usepackage[fontsize=10.5pt]{scrextend}
\usepackage[margin=22mm,headheight=15pt,headsep=9mm,footskip=11mm]{geometry}
\usepackage{fontspec}
\setmainfont{texgyrepagella}[Extension=.otf,UprightFont=*-regular,BoldFont=*-bold,ItalicFont=*-italic,BoldItalicFont=*-bolditalic]
\setsansfont{texgyreheros}[Extension=.otf,UprightFont=*-regular,BoldFont=*-bold,ItalicFont=*-italic,BoldItalicFont=*-bolditalic]
\setmonofont{texgyrecursor}[Extension=.otf,UprightFont=*-regular,BoldFont=*-bold,ItalicFont=*-italic,BoldItalicFont=*-bolditalic,Scale=MatchLowercase]
\usepackage{amsmath,amssymb}
\usepackage{unicode-math}
\setmathfont{texgyrepagella-math.otf}
\usepackage{xcolor,microtype,enumitem,needspace,fancyhdr}
\usepackage{bookmark}
\definecolor{ink}{HTML}{17344B}
\definecolor{accent}{HTML}{197D80}
\definecolor{muted}{HTML}{596773}
\hypersetup{colorlinks=true,linkcolor=accent,urlcolor=accent,pdftitle={PF-05 - Infinitesimal
rigidity detects unique size-two factors in the presence of
zeros},pdfauthor={Open Problems in Numerical Linear Algebra}}
\urlstyle{same}
\setlength{\parindent}{0pt}
\setlength{\parskip}{5pt plus 1pt minus 1pt}
\linespread{1.04}
\setlength{\emergencystretch}{3em}
\widowpenalty=10000
\clubpenalty=10000
\displaywidowpenalty=10000
\allowdisplaybreaks[1]
\setlist{nosep,leftmargin=1.5em}
\providecommand{\tightlist}{\setlength{\itemsep}{0pt}\setlength{\parskip}{0pt}}
\providecommand{\pandocbounded}[1]{#1}
\pagestyle{fancy}
\fancyhf{}
\fancyhead[L]{\sffamily\footnotesize\color{muted}OPEN PROBLEMS IN NUMERICAL LINEAR ALGEBRA}
\fancyhead[R]{\sffamily\bfseries\footnotesize\color{accent}PF-05}
\fancyfoot[L]{\parbox[b]{0.9\textwidth}{\sffamily\fontsize{8}{10}\selectfont\color{muted}Editorial ratings; a bounded search cannot certify that a problem remains unsolved.\\Literature check: 2026-09-11}}
\fancyfoot[R]{\sffamily\footnotesize\color{muted}\thepage}
\renewcommand{\headrulewidth}{0.3pt}
\renewcommand{\headrule}{\hbox to\headwidth{\color{accent}\leaders\hrule height \headrulewidth\hfill}}
\makeatletter
\renewcommand{\section}{\Needspace{5\baselineskip}\@startsection{section}{1}{0pt}{12pt plus 2pt minus 2pt}{4pt}{\sffamily\bfseries\large\color{ink}}}
\renewcommand{\subsection}{\Needspace{4\baselineskip}\@startsection{subsection}{2}{0pt}{9pt plus 1pt minus 1pt}{3pt}{\sffamily\bfseries\normalsize\color{ink}}}
\makeatother
\setcounter{secnumdepth}{0}
\begin{document}
{\sffamily\small\color{muted}Nonnegative and positive
factorizations\par}
\vspace{3pt}
{\raggedright\hyphenpenalty=10000\exhyphenpenalty=10000\sffamily\bfseries\fontsize{21}{24}\selectfont\color{ink}Infinitesimal
rigidity detects unique size-two factors in the presence of zeros\par}
\vspace{7pt}
{\sffamily\small\textbf{Difficulty:} challenging\quad\textbf{Importance:} interesting
to specialist\par}
{\sffamily\small\bfseries\color{accent}Status: Solved\par}
\vspace{4pt}
{\color{accent}\hrule height 0.8pt}
\vspace{6pt}
\textbf{Rating rationale:} Challenging because zeros change the feasible
infinitesimal motions and defeat the positive-entry argument; specialist
importance is a precise uniqueness test for size-two PSD factors.

\subsection{Resolution --- 2026-09-11}\label{resolution-2026-09-11}

\textbf{Author:} Matthew J. Colbrook, Department of Applied Mathematics
and Theoretical Physics, University of Cambridge. \textbf{Independent
Codex-agent review: PASS for the exact target.}

For every real size-two positive semidefinite factorization of an
ordinary-rank-three matrix, feasible straight-line infinitesimal
rigidity is equivalent to uniqueness up to congruence. The zero-entry
argument handles repeated or singular factors and zero rows and columns;
the cited positive-entry theorem covers the remaining case. The feasible
directions and equivalence group agree exactly with the canonical
definitions.

The complete target is resolved. Its former difficulty rating is
historical; the original statement, references and dated audits remain
below.

\textbf{Primary reference:}
\href{https://github.com/ajt60gaibb/OpenProblemsInNLA/blob/main/references/colbrook-factorization-2026-09-11/manuscripts/PF-05_rigidity_with_zeros.pdf}{complete
authored PDF},
\href{https://github.com/ajt60gaibb/OpenProblemsInNLA/blob/main/references/colbrook-factorization-2026-09-11/manuscripts/PF-05_rigidity_with_zeros.tex}{standalone
TeX}, \textbf{Theorem 1, with Theorem 8 for the zero-entry
construction}.
\href{https://github.com/ajt60gaibb/OpenProblemsInNLA/blob/main/references/colbrook-factorization-2026-09-11/verification/reviews/PF-05-review.md}{Independent
proof review} ·
\href{https://github.com/ajt60gaibb/OpenProblemsInNLA/blob/main/references/colbrook-factorization-2026-09-11/README.md}{Authorship
and submission record}. Verification is independent agent review, not
external human peer review or formal certification.

\subsection{Context and notation}\label{context-and-notation}

Write \(\mathbb S_+^k\) for the cone of real symmetric positive
semidefinite \(k\times k\) matrices. For an entrywise nonnegative matrix
\(M\in\mathbb R_+^{p\times q}\), its \textbf{real positive semidefinite
rank} is

\[
\operatorname{rank}_{\rm psd}(M)
=\min\{k\ge1:\ \exists A_1,\ldots,A_p,B_1,\ldots,B_q\in\mathbb S_+^k,\quad
M_{ij}=\operatorname{tr}(A_iB_j)\ \text{for all }i,j\}.
\]

This definition concerns a family of matrix factors indexed by the rows
and columns of \(M\).

\subsection{Problem statement}\label{problem-statement}

Let \(M\in\mathbb R_+^{p\times q}\) satisfy \(\operatorname{rank}(M)=3\)
and \(\operatorname{rank}_{\rm psd}(M)=2\), and fix any size-two
factorization \(M_{ij}=\operatorname{tr}(A_iB_j)\). Zero entries of
\(M\) are allowed.

Call a collection of real symmetric matrices \((E_i,F_j)\) a feasible
linear direction if

\[
\operatorname{tr}(E_iB_j)+\operatorname{tr}(A_iF_j)=0\quad\text{for every }i,j
\]

and some \(h>0\) satisfies \(A_i+tE_i\succeq0\), \(B_j+tF_j\succeq0\)
for all \(i,j\) and \(t\in[0,h)\). The trace constraint is imposed only
to first order; the perturbed factors need not represent \(M\) at
positive \(t\).

\newpage
\subsubsection{Question}\label{question}

Are the following conditions equivalent for every such factorization?

\begin{enumerate}
\def\labelenumi{\arabic{enumi}.}
\tightlist
\item
  Every feasible linear direction has the form \(E_i=dA_i\),
  \(F_j=-dB_j\) for one real scalar \(d\) common to all factors.
\item
  Every other size-two factorization \((\widetilde A_i,\widetilde B_j)\)
  of \(M\) satisfies \(\widetilde A_i=S^\mathsf T A_iS\) and
  \(\widetilde B_j=S^{-1}B_jS^{-\mathsf T}\) for one
  \(S\in GL(2,\mathbb R)\).
\end{enumerate}

This asks whether an infinitesimal test certifies uniqueness of the
optimal matrix factors up to changes of basis. Condition 1 is exactly
the source's 2-infinitesimal rigidity: for size two, the relevant
second-degree Taylor polynomials are the complete principal minors, and
Theorem 2.2 identifies the trivial directions as the common scalings
above.

\subsection{Reference and status}\label{reference-and-status}

K. Dawson, S. Hoşten, K. Kubjas, and L. Metsälampi,
\href{https://arxiv.org/html/2410.18891v2}{\emph{Uniqueness of size-2
positive semidefinite matrix factorizations}, v2}, 31 August 2026,
Definition 1, Remark 1, Theorem 2.2, and the conjecture immediately
after Theorem 5.2. Theorem 5.2 proves the assertion when every entry of
\(M\) is strictly positive. Lemma 19 identifies local and global
rigidity in the stated rank class, so omitting the equivalent local
condition does not weaken the source's remaining conjecture. Searches
for the title, PSD rigidity with zeros, and later work found no
resolution beyond this latest version.

\subsection{Status check --- 2026-09-10}\label{status-check-2026-09-10}

Rechecked
\href{https://arxiv.org/html/2410.18891v2}{Dawson--Hoşten--Kubjas--Metsälampi
v2, Theorem 2.2, Theorem 5.2 and its following conjecture}, and searched
for later rigidity results with zero entries. The August 31, 2026
revision proves equivalence for strictly positive M and conjectures its
extension to matrices with zeros. Those positive matrices are a proved
part of the displayed target. No resolution of the zero-entry remainder
was located.
\end{document}
